\documentclass[11pt,a4paper]{article}
\usepackage[margin=2.2cm]{geometry}
\usepackage{amsmath,amssymb,bm}
\usepackage{booktabs,array,longtable}
\usepackage{enumitem}
\usepackage{microtype}
\usepackage{xcolor}
\usepackage[colorlinks,linkcolor=red!50!black,citecolor=red!70!black,urlcolor=red!80!black]{hyperref}

% ---------------------------------------------------------------
% Notation (kept close to mr-hedin.tex and to Wang, Fang, and Li)
% ---------------------------------------------------------------
\newcommand{\bG}{\boldsymbol{G}}
\newcommand{\bSigma}{\boldsymbol{\Sigma}}
\newcommand{\bH}{\boldsymbol{H}}
\newcommand{\bM}{\boldsymbol{M}}
\newcommand{\bE}{\boldsymbol{E}}
\newcommand{\bF}{\boldsymbol{F}}
\newcommand{\bbH}{\mathbb{H}}
\newcommand{\bbT}{\mathbb{T}}
\newcommand{\ket}[1]{|#1\rangle}
\newcommand{\bra}[1]{\langle#1|}
\newcommand{\braket}[1]{\langle#1\rangle}
\newcommand{\GW}{$GW$}
\newcommand{\MRGW}{MR-$GW$}
\newcommand{\WFL}{WFL}
\newcommand{\Pre}{Preliminary section}
\newcommand{\DHN}{Dyall--Hedin notes}
\newcommand{\Sc}[1]{\textsf{[#1]}}

\title{Comparison of the preliminary section ``Multireference Green's Function methods''
with the \MRGW{} paper of Wang, Fang, and Li (arXiv:2604.16013) and with the
Dyall--Hedin notes in \texttt{mr-hedin.tex}}
\author{}
\date{September 22, 2026}

\begin{document}
\maketitle

\section{Scope and the three texts compared}

Three texts are compared.

\begin{description}[style=nextline,leftmargin=1.5em]
\item[(A) The \emph{\Pre}.] The section ``Multireference Green's Function
  methods'' at the beginning of \texttt{mr-hedin.tex}.  It partitions the exact
  self-energy into an ``initial'' part $\bSigma^{\text{init}}$ and a
  ``correlation'' part $\bSigma^{\text{corr}}$, takes the initial part from an
  active-space calculation, and constructs the corresponding initial Green's
  function $\bG_{\text{init}}(\omega)$ directly as the resolvent of a block
  matrix in the space spanned by the inactive spin orbitals and the exact
  $N\pm1$ eigenstates of the active-space Hamiltonian.  The inactive and
  active blocks are coupled by a matrix $\bH_{\mathrm{i,a}}$ of double-commutator
  (superoperator) matrix elements.  Block inversion is then shown to produce a
  factorized form that has the same appearance as the generalized Dyson equation
  of the \MRGW{} paper.
\item[(B) \WFL{}.] Y.~Wang, W.-H.~Fang, and Z.~Li, ``A Diagrammatic Route to
  Multi-Reference $GW$ for Strongly Correlated Molecules'', arXiv:2604.16013v2
  (3~Sep~2026)~\cite{WFL2026}.  The paper builds a diagrammatic many-body
  perturbation theory on the interacting Dyall Hamiltonian
  $\hat H_0^{\text{Dyall}}$, organizes the expansion with Hall's generalized
  Dyson equation (four self-energy blocks $\Sigma^{11},\Sigma^{12},\Sigma^{21},\Sigma^{22}$),
  derives the first-order blocks (their Appendix~B and Sec.~S1), and defines a
  practical one-shot \MRGW{} approximation in which only $\Sigma^{11}$ is kept,
  with the residual interaction screened at the MR-RPA level.  The mixed blocks
  $\Sigma^{12/21}$ are deliberately dropped because their isolated inclusion
  breaks positive semidefiniteness (their Appendix~C, Fig.~6).  Benchmarks on
  Be, stretched H$_2$, and O$_3$ are reported.
\item[(C) The \emph{\DHN}.] The remainder of \texttt{mr-hedin.tex} (Objectives~I
  and~II).  Objective~I embeds the non-Gaussian Dyall reference into an enlarged
  source space following Brouder~\cite{Brouder2009}, recovers Hall's
  generalized Dyson equation and the four first-order blocks of \WFL{}, and
  identifies \MRGW{} as a channel-selective truncation of an exact
  Dyall--Hedin hierarchy.  Objective~II restores the mixed blocks with a
  statically screened kernel $v_R-[W_D(0)]^x$ and embeds them, together with
  the dynamic \MRGW{} self-energy, in a common Hermitian pole-space matrix that
  is causal, positive semidefinite (PSD), and moment conserving.
\end{description}

The comparison is organized as follows.  Section~\ref{sec:dictionary} gives a
dictionary of notation.  Section~\ref{sec:similarities} lists the points on
which the three texts agree, including several places where the
\Pre{} anticipates constructions that \WFL{} call for and the
\DHN{} carry out.  Section~\ref{sec:differences} lists the differences, some of
which are differences of formalism and some of which are genuine
discrepancies in content.  Section~\ref{sec:synthesis} describes how the
\Pre{} fits into the rest of the notes, and Sec.~\ref{sec:revisions} collects
the concrete points in the \Pre{} that the comparison suggests revising.

\section{Dictionary of notation}
\label{sec:dictionary}

The same physical objects appear under different names in the three texts.
Table~\ref{tab:dictionary} lists the correspondences.  Throughout, $P,Q$ label
inactive spin orbitals (core $i,j$ and virtual $a,b$), $x,y,z,w$ label active
spin orbitals, $\mu$ labels eigenstates of the active-space Hamiltonian, and
$\alpha=(\mu,\pm)$ labels an active electron-attached or electron-removed pole.

\begin{table}[h!]
\centering
\small
\begin{tabular}{>{\raggedright\arraybackslash}p{3.1cm}>{\raggedright\arraybackslash}p{4.2cm}>{\raggedright\arraybackslash}p{3.8cm}>{\raggedright\arraybackslash}p{3.3cm}}
\toprule
\Pre{} & \WFL{} & \DHN{} & Meaning \\
\midrule
$\bG_0(\omega)$ & $g_0$ [Eq.~(10)] & $g_0$ & non-interacting (Fock-type) propagator \\[1mm]
$\bSigma^{\text{init}}$ & $\Sigma_0[g_0,\hat V^A]=g_0^{-1}-G_0^{-1}$ [Eq.~(10)] & $\Sigma_D=g_0^{-1}-G_D^{-1}$ & self-energy carrying the nonperturbative active-space correlation \\[1mm]
$\bG_{\text{init}}$ with $\bH_{\mathrm{i,a}}=0$ & $\mathbf G_0=\operatorname{diag}(G_0^I,G_0^A)$ [Eq.~(3)] & $G_D$ & Dyall Green's function \\[1mm]
$\bG_{\text{init}}$ with $\bH_{\mathrm{i,a}}\neq0$ & no direct counterpart; closest object is MR-GF1, Eqs.~(S18)--(S20) & $\bbT(z-\bbH)^{-1}\bbT^\dagger$ with $A=0$ and the bare kernel $\bar v_R$ & Hermitian embedding of the inactive--active coupling \\[1mm]
$\bSigma^{\text{corr}}$ & $\Sigma^{\text{MR-}GW}$ as an approximation to $\Sigma^{11}$ [Eqs.~(9), (S41)] & all of $\Sigma^{11},\Sigma^{12},\Sigma^{21},\Sigma^{22}$ generated by $\hat V_R$ & everything not contained in the reference \\[1mm]
$\ket{\Phi_{\mathrm i}}\otimes\ket{\Psi_{\mathrm a}}$ & $\ket{\Phi^I_{\mu_I}}\ket{\Xi^A_{\mu_A}}$ & Dyall eigenstates (inactive Gaussian $\otimes$ active correlated) & product structure of the reference \\[1mm]
$\bF^{\text{HF}}_{\mathrm{i,i}}$ & $\epsilon_P$ from the Dyall Fock operator, Eq.~(14), which contains the active one-body density & $\kappa_P=\epsilon_P$ & inactive one-particle energies \\[1mm]
$\bE^{N\pm1}$ & $\omega^A_{\mu,N_a\pm1}$ [Eq.~(S17)] & signed $\kappa_{\mu\pm}$ in $K_D$ & active charged excitation energies \\[1mm]
$\bM^{N\pm1}_{p\mu}$ & $D^{[\pm1]}_{\mu,x}$ [Eqs.~(S49)--(S50)] & $d_\alpha$, columns of $T_D$ & active Dyson amplitudes \\[1mm]
$\bH_{\mathrm{i,a}}$ & $[\Sigma_1^{11}]_{Px}D^{[\pm1]}_{\mu,x}$ plus the residues of $\Sigma_1^{12,IA}$ [Eqs.~(S9)--(S12), (S15)] & $[T_D^\dagger\Sigma_{\rm stat}T_D+K_B^{(1)}]_{P\alpha}$ & bare first-order inactive--active coupling \\[1mm]
$\bSigma_{\mathrm{i,a}}(\omega)$ & $\Sigma_1^{12,IA}(\omega)$ [Eq.~(S15)], plus the static $\Sigma_1^{11}$ piece & $\Sigma^{12,R}_{\rm scr}$ in the bare limit & right mixed block \\[1mm]
$\bSigma_{\mathrm{a,i}}(\omega)$ & $\Sigma_1^{21,AI}(\omega)$ [Eq.~(S16)] & $\Sigma^{21,R}_{\rm scr}$ in the bare limit & left mixed block \\[1mm]
$\bSigma_{\mathrm{i,i}}(\omega)$ & absent from \MRGW; would be a second-order contribution to $\Sigma^{11}$ & second-order term of the $K_B$ resummation & inactive self-energy with active $N\pm1$ poles \\[1mm]
$\tilde\bSigma_{\mathrm{i,i}}=\bG_{\mathrm i}-\bG_{\mathrm{i},0}$ & identified in (A) with the $\Sigma^{22}$ slot of Eq.~(5); see Sec.~\ref{sec:differences}, item~D3 & --- & inactive dressing, not a self-energy \\[1mm]
$(a|b)=\braket{\Psi_0|[a^\dagger,b]|\Psi_0}$, killer condition & --- & retarded representation; $T_DT_D^\dagger=I$ & superoperator metric; zeroth-moment sum rule \\[1mm]
$\{\hat C^\dagger\}=\{a_i,a_a\}$ & --- & inactive pole vectors $t_P=e_P$ & primary (inactive) operator manifold \\
\bottomrule
\end{tabular}
\caption{Correspondence of notation.  Equation numbers refer to
\WFL{} v2 (main text, End Matter, and Supplemental Material).}
\label{tab:dictionary}
\end{table}

\section{Similarities}
\label{sec:similarities}

\paragraph{S1. Same partition of the problem and the same ``reference plus
correction'' split of the self-energy.}
The \Pre{} writes $\bSigma=\bSigma^{\text{init}}+\bSigma^{\text{corr}}$ and
defines $\bG_{\text{init}}=(\bG_0^{-1}-\bSigma^{\text{init}})^{-1}$ so that
$\bG=\bG_{\text{init}}+\bG_{\text{init}}\bSigma^{\text{corr}}\bG$.  This is,
term for term, the ``connection to standard self-energy'' of \WFL{}, Eqs.~(10)
and~(11): a non-interacting $g_0$ is connected to the interacting Dyall
Green's function $G_0$ by $\Sigma_0[g_0,\hat V^A]=g_0^{-1}-G_0^{-1}$, and the
full $G$ then satisfies
$G=g_0+g_0\,(\Sigma_0+\Sigma^{\text{MR-}GW})\,G$.  The \DHN{} restate the same
relation in the subsection ``Why the construction is internally consistent''
with $\Sigma_D=g_0^{-1}-G_D^{-1}$.  In all three texts, the active-space
two-electron interaction lives entirely in the reference and only the residual
interaction $\hat V_R=\hat H-\hat H_D$ is treated by the correction, so that
no double-counting correction is needed.  This is the multireference
perturbation-theory (CASPT2/NEVPT2) philosophy transferred to the one-particle
propagator, and it is shared by (A), (B), and (C).

\paragraph{S2. Same factorized reference and block-diagonal reference
propagator.}
The product form $\ket{\Psi_0}=\ket{\Phi_{\mathrm i}}\otimes\ket{\Psi_{\mathrm a}}$
assumed in the \Pre{} is exactly the factorization
$\ket{\Phi_\mu}=\ket{\Phi^I_{\mu_I}}\ket{\Xi^A_{\mu_A}}$ that \WFL{} obtain
from the additive structure of the Dyall Hamiltonian, and it is what makes the
inactive sector Gaussian in the \DHN{}.  The consequence noted in the
\Pre{}, ``there are no Dyson amplitudes that couple the active and inactive
spaces directly'', is \WFL{}'s Eq.~(3), $\mathbf G_0=\operatorname{diag}(G_0^I,G_0^A)$,
and is the statement in the \DHN{} that connected Dyall cumulants of order
$\ge2$ carry only active indices (``Specialization to the Dyall
Hamiltonian'').  The active Dyson amplitudes
$M_{p\mu}=\braket{\Psi_0|a_p|\Psi_\mu^{N+1}}$ of the \Pre{} are the transition
densities $D^{[+1]}_{\mu,x}$ and $D^{[-1]}_{\mu,x}$ of \WFL{} [Eqs.~(S49)--(S50)]
and the pole vectors $d_\alpha$ of the \DHN{}.

\paragraph{S3. Same block structure of the inactive--active coupling.}
Block inversion in the \Pre{} gives
\begin{equation}
  \bG_{\text{init}}
  =
  \begin{pmatrix}
    \bG_{\mathrm i} & \bG_{\mathrm i}\bSigma_{\mathrm{i,a}} \\
    \bSigma_{\mathrm{a,i}}\bG_{\mathrm i} & \bG_{\mathrm A}+\bSigma_{\mathrm{a,i}}\bG_{\mathrm i}\bSigma_{\mathrm{i,a}}
  \end{pmatrix},
  \qquad
  \bSigma_{\mathrm{i,a}}(\omega)=-\bH_{\mathrm{i,a}}\,(\omega-\bE^{N\pm1})^{-1}\bM^\dagger .
  \label{eq:pre-Ginit2}
\end{equation}
This is precisely the structure of \WFL{}'s Eq.~(S20),
\begin{equation}
  \mathbf M
  =
  \begin{pmatrix}
    G_0^I & G_0^I\Sigma_1^{12,IA} \\
    \Sigma_1^{21,AI}G_0^I & \Sigma_1^{21,AI}G_0^I\Sigma_1^{12,IA}+G_0^A
  \end{pmatrix},
  \label{eq:wfl-S20}
\end{equation}
which follows from Hall's generalized Dyson equation once one uses the
triangular block structure $\Sigma_1^{12}=\bigl(\begin{smallmatrix}0&\Sigma^{12,IA}_1\\0&0\end{smallmatrix}\bigr)$,
$\Sigma_1^{21}=\bigl(\begin{smallmatrix}0&0\\ \Sigma^{21,AI}_1&0\end{smallmatrix}\bigr)$
and $\Sigma_1^{22}=0$ [\WFL{} Eq.~(S14); \DHN{} ``Triangular structure of the
mixed blocks'' and ``Vanishing of $\Sigma_1^{22}$''].  The \Pre{} recognizes
this (``Eq.~(Ginit3) has the same structure as the initial Dyson equation in
the \MRGW{} paper''), and the recognition is correct at the level of the block
pattern.  Moreover, the pole content agrees: $\bSigma_{\mathrm{i,a}}$ has
simple poles at the active $N\pm1$ excitation energies, exactly as
$\Sigma_1^{12,IA}$ in \WFL{}'s Eq.~(S15) and as the retarded mixed blocks of
the \DHN{} (``Retarded representation of the screened mixed blocks'').

\paragraph{S4. The coupling matrix $\bH_{\mathrm{i,a}}$ is the first-order
residual coupling and contains the same three-active-operator amplitude as the
mixed Hall blocks.}
The double-commutator element
$\bH_{P\mu}=\braket{\Psi_0|[a_P,[\hat H,\hat R^{N\pm1}_\mu]]|\Psi_0}$ of the
\Pre{} vanishes for the Dyall part of $\hat H$ (the inactive operator $a_P$
anticommutes with any active operator, and $[\hat H_D,\hat R_\mu]$ is again an
active operator) and is therefore of first order in $\hat V_R$.  Evaluating it
for the electron-attachment sector with $\hat V_R=u_{pq}\hat p^\dagger\hat q+\tfrac14\bar v_{pr,qs}\hat p^\dagger\hat q^\dagger\hat s\hat r$
and a virtual orbital $P=a$ gives
\begin{equation}
  \bH_{a\mu}
  =
  \bigl(u_{az}+\bar v_{az,kk}\bigr)\,D^{[+1]}_{\mu,z}
  +
  \tfrac12\,\bar v_{az,yw}\,
  \braket{\Xi^A_{0,N_a}|\hat y^\dagger\hat w\hat z|\Xi^A_{\mu,N_a+1}} ,
  \label{eq:H-ia-explicit}
\end{equation}
where $k$ runs over core orbitals.  Adding and subtracting the mean-field
piece $\bar v_{az,yw}\gamma_{yw}D^{[+1]}_{\mu,z}$ shows that
\begin{equation}
  \bH_{a\mu}
  =
  \underbrace{[\Sigma_1^{11}]_{az}\,D^{[+1]}_{\mu,z}}_{\text{\WFL{} Eq.~(S12)}}
  +
  \underbrace{\Bigl[\tfrac12\bar v_{az,yw}\braket{\Xi_0|\hat y^\dagger\hat w\hat z|\Xi_\mu^{+}}
  -\bar v_{az,yw}\gamma_{yw}D^{[+1]}_{\mu,z}\Bigr]}_{\text{residue of \WFL{} Eq.~(S15); }\bar v_R\,\mathcal C^{D,+}_{\mu;z,yw}\text{ in the \DHN}} .
  \label{eq:H-ia-split}
\end{equation}
The second bracket is the ``connected'' amplitude that \WFL{} identify as the
residue of $\Sigma_1^{12,IA}$ and that the \DHN{} write as
$\bar v_R\,\mathcal C^{D,+}_{\mu;z,yw}$ (``Screened mixed residues'').  The
\DHN{} make the same observation in the opposite direction in the subsection
``Same-sector pole pairs'': in the bare limit the mixed residue combined with the
off-diagonal $\Sigma_1^{11}$ contribution ``reconstructs the first-order
residual-Hamiltonian coupling between the two charged Dyall states''.  That
first-order coupling is exactly $\bH_{\mathrm{i,a}}$.  The remark in the \Pre{}
that only a limited number of excitation manifolds in $\hat R^{N\pm1}$ need to
be considered because $\hat H$ is at most two-body is the same fact that lets
the \DHN{} replace the explicit three-operator tensor by a contracted source
vector $\hat Q_P\ket{\Xi_0}$ (``Avoiding the explicit three-operator transition
tensor'').

\paragraph{S5. The \Pre{} is a Hermitian (superoperator) embedding, and it is
the bare-coupling, no-bath limit of the common Hermitian matrix of
Objective~II.}
This is the most important point of contact.  The \Pre{} never solves Hall's
equation; it defines $\bG_{\text{init}}$ as the resolvent of a Hermitian block
matrix,
\begin{equation}
  \bG_{\text{init}}(\omega)
  =
  \begin{pmatrix}\mathbf 1&0\\0&\bM\end{pmatrix}
  \begin{pmatrix}
    \omega\mathbf 1-\bF_{\mathrm{i,i}} & -\bH_{\mathrm{i,a}}\\
    -\bH_{\mathrm{i,a}}^\dagger & \omega\mathbf 1-\bE^{N\pm1}
  \end{pmatrix}^{-1}
  \begin{pmatrix}\mathbf 1&0\\0&\bM^\dagger\end{pmatrix},
  \label{eq:pre-Ginit}
\end{equation}
(up to the sign convention of the off-diagonal block, cf.\ item~D6).  With the
identifications $t_P=e_P$, $d_\alpha=\bM_{\cdot\alpha}$, $K_D=\operatorname{diag}(\epsilon_P,\kappa_\alpha)$,
$T_D=(\mathbf 1\;\;\bM)$, this is
\begin{equation}
  \bG_{\text{init}}(z)
  =
  T_D\bigl(zI-K_{\rm top}\bigr)^{-1}T_D^\dagger ,
  \qquad
  K_{\rm top}=K_D+T_D^\dagger\Sigma_{\rm stat}T_D+K_B^{(1)},
  \label{eq:pre-as-Ktop}
\end{equation}
i.e.\ the common Hermitian matrix $\bbH_{\text{mix-MR-}GW}$ of the \DHN{}
(``Common Hermitian matrix'') with the \MRGW{} bath switched off ($A=0$) and
the mixed kernel taken at its bare value $\mathcal I_D^{(0)}=\bar v_R$
(``Hierarchy of static mixed kernels'').  Consequently, the \Pre{} construction
automatically has the three properties that Objective~II sets out to secure:
it is causal, its spectral function is PSD, and it satisfies the zeroth-moment
sum rule $\int A(\omega)\,d\omega=I$ because $T_DT_D^\dagger=\mathbf 1$ (the
active part $\bM\bM^\dagger=\mathbf 1_{\rm act}$ is the completeness of the
$N_a\pm1$ active eigenstates).  Expanding the resolvent in $K_B$ shows also
that the inactive self-energy of the \Pre{},
$\bSigma_{\mathrm{i,i}}(\omega)=\bH_{\mathrm{i,a}}(\omega-\bE)^{-1}\bH_{\mathrm{i,a}}^\dagger$,
is the second-order term $R_DK_BR_DK_BR_D$ of the $K_B$ resummation that
Objective~II calls a ``chosen Hermitian resummation''.  The \Pre{} thus
provides an independent derivation of the same higher-order completion.

\paragraph{S6. The ``direct determination without the Dyson equation'' is the
linearization lemma of Objective~II.}
The \Pre{} remarks that $\bG_{\text{init}}$ can be obtained directly, without
the Dyson equation.  The \DHN{} prove the identity
$T[R^{-1}-T^\dagger XT]^{-1}T^\dagger=[G_0^{-1}-X]^{-1}$ for rectangular $T$
(``A useful linearization identity'', a Woodbury-type identity), which is
precisely the statement that the resolvent route of the \Pre{} and the
Dyson-equation route are interchangeable, and which also justifies folding the
\MRGW{} self-energy into the same Hermitian matrix.

\paragraph{S7. The same two missing ingredients are identified, and the same
type of PSD completion is called for.}
The \Pre{} ends by listing what $\bSigma^{\text{corr}}$ still has to provide:
(i) correlation within the inactive space and (ii) inactive--active couplings
beyond the operator manifold $\{a_P\}\cup\{\hat R^{N\pm1}_\mu\}$.  \WFL{}
provide (i) and part of (ii) through the four MR-RPA screening channels
[their Fig.~1(d): core$\to$virtual, core$\to$active, active$\to$virtual,
active$\to$active], and the \DHN{} classify what remains as the channel-selective
truncation of Objective~I.  Conversely, \WFL{}'s Appendix~C states that a
consistent treatment of $\Sigma^{12/21}$ ``will require an order-consistent
higher-order completion or resummation'' and points to PSD-preserving ADC-type
reformulations of vertex-corrected \GW{}~\cite{MarieTolleLoos2026}.  The
\Pre{} (with bare couplings) and Objective~II (with screened couplings and the
\MRGW{} bath) are both completions of exactly that kind.

\section{Differences}
\label{sec:differences}

\paragraph{D1. Formalism: superoperator versus diagrammatic versus functional.}
The \Pre{} is written in the language of electron-propagator (superoperator)
theory familiar from EOM-IP/EA, ADC, and MR-ADC~\cite{GoscinskiLukman1970,Ortiz2013,Schirmer1983,Sokolov2018,ChatterjeeSokolov2019}:
a binary product $(a|b)=\braket{\Psi_0|[a^\dagger,b]_+|\Psi_0}$, a killer
condition, and a superoperator Hamiltonian matrix
$\braket{\Psi_0|[\hat C,[\hat H,\hat R]]|\Psi_0}$.  \WFL{} use the
Gell-Mann--Low expansion in the interaction picture generated by the
\emph{interacting} $\hat H_0$, followed by the cumulant decomposition of the
two- and three-body Dyall Green's functions [their Eqs.~(S1)--(S7)].  The
\DHN{} use Schwinger's functional-derivative technique with Brouder's
enlarged source space, in which the cumulants become auxiliary vertices and an
ordinary matrix Dyson equation holds in the enlarged space.

The practical consequence concerns Hermiticity and time ordering.  The
anticommutator metric of the \Pre{} produces the \emph{retarded} (combined
IP/EA) propagator, so that the $N+1$ and $N-1$ poles sit in one Hermitian
matrix with a common $+i0^+$ prescription and the coupling satisfies
$\bH_{\mathrm{a,i}}=\bH_{\mathrm{i,a}}^\dagger$ by construction.  The blocks
$\Sigma_1^{12}$ and $\Sigma_1^{21}$ of \WFL{} are \emph{time-ordered} objects
[Eqs.~(S15)--(S16)]: the attachment and removal terms carry opposite
$\pm i0^+$ prescriptions and the two blocks are not simply Hermitian
conjugates.  The \DHN{} must therefore convert to the retarded representation
before building $K_B$ (``Retarded representation of the screened mixed blocks'',
``Cross-sector pole pairs'').  The \Pre{} reaches the same retarded, Hermitian
$K_B$ in one step, which is a genuine advantage of its formulation.  The
red-marked warning in the \Pre{} about ``some problem with the sign convention''
corresponds precisely to the bookkeeping that \WFL{} fix in Eqs.~(20)--(23)
and (S5)--(S7) and that the \DHN{} fix through the source-transfer identity and
the subsection ``Origin of the numerical prefactors'' (the factors $-i/2$ and
the connected subtraction $-\bar v\gamma G^A$).  The explicit evaluation in
Eq.~\eqref{eq:H-ia-split} above shows that the prefactor $\tfrac12$ and the
mean-field subtraction come out of the double commutator automatically once
the antisymmetrized interaction is used.

One qualification is needed, and it has been verified numerically (see
\texttt{numerics/run\_checks.py}).  With the zeroth-order $\ket{\Psi_0}$ in
the binary product, the double commutator of the \Pre{} produces only the
\emph{same-sector} couplings: a virtual orbital $a$ with the $N_a+1$ states
and a core orbital $i$ with the $N_a-1$ states.  For these pairs it equals the
Hamiltonian matrix element between the two orthonormal charged Dyall states,
$\bH_{a\mu}=\braket{a_a^\dagger\Psi_0|\hat H|\Phi_I\,\Xi_\mu^{N_a+1}}$ and
$\bH_{i\mu}=-\braket{a_i\Psi_0|\hat H|\Phi_I\,\Xi_\mu^{N_a-1}}$, because
$\bra{\Psi_0}a_i=0$ for a core orbital and $a_a\ket{\Psi_0}=0$ for a virtual
one.  The \emph{cross-sector} entries of $K_B$ in the \DHN{} (core $i$ with
$N_a+1$ states, virtual $a$ with $N_a-1$ states) are not matrix elements of
any $N\pm1$ Hamiltonian.  They encode the first-order change of the Dyson
amplitudes caused by the admixture of core-hole$\,\otimes(N_a+1)$ and
virtual-particle$\,\otimes(N_a-1)$ configurations into the ground state,
i.e.\ residual-interaction ground-state correlation, which in the
superoperator formalism enters only through a correlated reference, exactly
as in ADC.  Both kinds of entries are required to reproduce the exact
first-order Green's function: with the same-sector entries alone the
finite-difference derivative $dG/d\lambda$ of the full-CI Green's function of
$\hat H_D+\lambda\hat V_R$ is not recovered, with both it is recovered to the
finite-difference accuracy.

\paragraph{D2. What is resummed: bare all-order inactive--active mixing (A)
versus screened $\Sigma^{11}$ only (B) versus both (C).}
Table~\ref{tab:content} summarizes the content of the three constructions.
The \Pre{} treats the inactive--active coupling to all orders through the
$2\times2$ block resolvent, but the coupling is \emph{bare}, there is no
screening of the residual interaction anywhere, and the inactive--inactive
sector contains no dynamic correlation apart from $\bSigma_{\mathrm{i,i}}$
(which has active poles only).  \WFL{} keep only $\Sigma^{11}$, but inside it
they screen the residual interaction dynamically with the MR-RPA response and
thereby include inactive-space ring correlation and the bath of
$(N\pm1)\otimes$(neutral MR-RPA excitation) intermediate states; the mixed
blocks are set to zero.  Objective~II of the \DHN{} contains both: the mixed
couplings, statically screened, and the dynamic \MRGW{} bath, in one Hermitian
matrix.

\begin{table}[h!]
\centering
\small
\begin{tabular}{>{\raggedright\arraybackslash}p{3.9cm}>{\raggedright\arraybackslash}p{3.3cm}>{\raggedright\arraybackslash}p{3.3cm}>{\raggedright\arraybackslash}p{3.9cm}}
\toprule
Ingredient & \Pre{} (A) & \WFL{} \MRGW{} (B) & \DHN{} Objective~II (C)\\
\midrule
Exact active $N\pm1$ poles and amplitudes & yes & yes & yes\\
Static one-body residual term $\Sigma_1^{11}$ & inactive--active part only (in $\bH_{\mathrm{i,a}}$); core--virtual part absent & yes & yes\\
Inactive--active mixed coupling ($\Sigma^{12/21}$) & yes, bare, all orders & no & yes, statically screened, Hermitian resummation\\
Screened residual interaction $W_D$ & no & MR-RPA, dynamic & MR-RPA, dynamic in $\Sigma^{11}$, static in the mixed kernel\\
Dynamic \GW{}-type self-energy in $\Sigma^{11}$ & no & yes & yes\\
$\Sigma^{22}$ & no ($\tilde\bSigma_{\mathrm{i,i}}$ is not $\Sigma^{22}$; see D3) & $\Sigma_1^{22}=0$; higher orders dropped & $\Sigma_1^{22}=0$; higher orders dropped\\
Causality / PSD / zeroth moment & yes, by construction & yes for stable MR-RPA & yes, by construction\\
First-order exactness of $G$ & yes (bare limit) & no (mixed blocks missing) & yes\\
Exact first spectral moment through $\mathcal O(\hat V_R)$ & yes & no & yes\\
Numerical results & none & Be, H$_2$, O$_3$, H$_4$ & none\\
\bottomrule
\end{tabular}
\caption{Content of the three constructions.}
\label{tab:content}
\end{table}

\paragraph{D3. The identification $\tilde\bSigma_{\mathrm{i,i}}\leftrightarrow\Sigma^{22}$
is only structural and is not correct in the Hall--Brouder sense.}
The \Pre{} writes
$\bG_{\text{init}}=(\mathbf 1-\bSigma_{\mathrm{a,i}})^{-1}(\bG_{\mathrm{i},0}+\tilde\bSigma_{\mathrm{i,i}})(\mathbf 1-\bSigma_{\mathrm{i,a}})^{-1}$
and matches it to Hall's
$\mathbf M=(I-\Sigma^{21})^{-1}(\mathbf G_0+\Sigma^{22})(I-\Sigma^{12})^{-1}$,
which places $\tilde\bSigma_{\mathrm{i,i}}=\bG_{\mathrm i}-\bG_{\mathrm i,0}$
in the slot of $\Sigma^{22}$.  In the theory of \WFL{} and of the \DHN{},
$\Sigma^{22}$ is built from the auxiliary (cumulant) sector; its external
indices are those of a connected active-space cumulant, so it is supported on
the active block, and for the Dyall partition it vanishes at first order
[\WFL{} Appendix~C; \DHN{} ``Vanishing of $\Sigma_1^{22}$'']. The object
$\tilde\bSigma_{\mathrm{i,i}}$ is instead (a) supported on the \emph{inactive}
block, (b) of second order in $\hat V_R$, and (c) a Green's-function-like
quantity ($\bG_{\mathrm i,0}\bSigma_{\mathrm{i,i}}\bG_{\mathrm i}$) rather than
a self-energy.  In Hall's classification the underlying
$\bSigma_{\mathrm{i,i}}=\bH_{\mathrm{i,a}}(\omega-\bE)^{-1}\bH_{\mathrm{i,a}}^\dagger$
is a second-order contribution to $\Sigma^{11}$: it is the 2h1p/2p1h-type
self-energy in which the three active lines are replaced by the exact active
three-body propagator.  In the language of Objective~II it is the second-order
term of the $K_B$ resummation (cf.\ S5).  The block pattern of
Eq.~\eqref{eq:pre-Ginit2} therefore coincides with Eq.~\eqref{eq:wfl-S20}, but
the mapping of the middle factor onto $\Sigma^{22}$ should be dropped.

\paragraph{D4. Same block pattern, different analytic structure: MR-GF1 can
violate positivity, $\bG_{\text{init}}$ cannot.}
\WFL{} define MR-GF1 by inserting the bare first-order blocks into Hall's
equation, Eqs.~(S18)--(S20), and note that the term
$\Sigma_1^{21,AI}G_0^I\Sigma_1^{12,IA}$ produces incorrect analytic structure
because both mixed blocks have the poles of $G_0^A$.  The product then has
\emph{double} poles at the active $N\pm1$ energies, the residues are no longer
PSD, and the spectral function of H$_4$ becomes negative in places (their
Fig.~6).  This is the reason they omit $\Sigma^{12/21}$ from \MRGW{}.  The
corresponding term in the \Pre{}, $\bSigma_{\mathrm{a,i}}\bG_{\mathrm i}\bSigma_{\mathrm{i,a}}$,
contains the \emph{dressed} $\bG_{\mathrm i}=[\bG_{\mathrm i,0}^{-1}-\bSigma_{\mathrm{i,i}}]^{-1}$.
Near an active pole $\omega\to\omega_\mu$ one has
$\bSigma_{\mathrm{i,i}}\sim \bH_{\cdot\mu}\bH_{\cdot\mu}^\dagger/(\omega-\omega_\mu)$,
hence $\bG_{\mathrm i}\sim(\omega-\omega_\mu)$, and the apparent double pole
of $\bSigma_{\mathrm{a,i}}\bG_{\mathrm i}\bSigma_{\mathrm{i,a}}$ collapses to a
simple pole with a PSD residue, as it must for the resolvent of a Hermitian
matrix.  The statement in the \Pre{} that its result ``has the same structure
as the initial Dyson equation in the \MRGW{} paper'' is thus true of the block
pattern but not of the analytic content: the \Pre{} construction is the
cured version of MR-GF1, and the cure is exactly the Hermitian embedding that
Objective~II adopts and that \WFL{} ask for.  (The dressing of $\bG_{\mathrm i}$
is not an optional refinement; it is what removes the pathology.)

\paragraph{D5. Reference one-particle energies and the core--virtual block.}
The \Pre{} uses a Hartree--Fock Fock matrix $\bF^{\text{HF}}_{\mathrm{i,i}}$
for the inactive block.  \WFL{} and the \DHN{} use the Dyall Fock operator,
$F_{PQ}=h_{PQ}+\braket{Pk||Qk}+\braket{Px||Qy}\braket{\hat x^\dagger\hat y}$
[\WFL{} Eq.~(14)], which contains the active one-body density.  Only with this
choice is $\bG_{\text{init}}|_{\bH=0}$ the Green's function of $\hat H_D$, is
$[\Sigma_1^{11}]_{ij}=[\Sigma_1^{11}]_{ab}=[\Sigma_1^{11}]_{xy}=0$ [\WFL{}
Eqs.~(S8), (S11), (S13)], and does the residual interaction contain no
active--active one-body piece.  Using $\bF^{\text{HF}}$ instead shifts
one-body terms between $\hat H_D$ and $\hat V_R$ and changes the meaning of
$\bSigma^{\text{corr}}$.  A second, smaller point: the inactive block in the
\Pre{} is taken diagonal, so the core--virtual coupling
$[\Sigma_1^{11}]_{aj}=h_{aj}+\bar v_{aj,pq}\gamma_{pq}$ [\WFL{} Eq.~(S9)] is
absent.  For CASSCF-optimized orbitals this element vanishes by the
core--virtual stationarity condition, but for CASCI on, e.g., RHF orbitals
(which \WFL{} also use; their Table~S2) it does not.  Objective~II keeps it in
$T_D^\dagger\Sigma_{\rm stat}T_D$.

\paragraph{D6. Signs of the removal poles and the placement of $\bM$ versus $\bM^\dagger$.}
The \Pre{} writes a single diagonal block $\bE^{N\pm1}$ without stating that,
in a retarded (superoperator) matrix, the removal energies enter with a minus
sign, $\kappa_{\mu-}=-(E_\mu^{N-1}-E_0^N)$, and it writes
$\bSigma_{\mathrm{i,a}}=-\bH_{\mathrm{i,a}}(\omega-\bE)^{-1}\bM$ where the block
inverse of Eq.~\eqref{eq:pre-Ginit} gives $\bM^\dagger$.  The \DHN{} make both
points explicit (signed pole energies in $K_D$; the conjugation pattern
$[K_B]_{P\alpha}=B_{P\alpha}$, $[K_B]_{\alpha P}=B_{P\alpha}^*$).  The
off-diagonal sign in Eq.~\eqref{eq:pre-Ginit} is immaterial for
$\bG_{\text{init}}$ as long as the two off-diagonal blocks are Hermitian
conjugates, but it does fix the sign of $\bSigma_{\mathrm{i,a}}$ relative to
$\Sigma_1^{12,IA}$ of \WFL{}, which is what the red warning is about.

\paragraph{D7. Route to systematic improvement.}
The three texts point to different, complementary ways of going beyond their
respective starting points.  The \Pre{} speaks of couplings ``beyond the
operator manifold'' defined by $\{a_P\}\cup\{\hat R^{N\pm1}_\mu\}$, which is
the EOM/ADC route: enlarge the manifold by 2h1p/2p1h-type operators with
inactive indices, as in MR-ADC~\cite{Sokolov2018,ChatterjeeSokolov2019}.
\WFL{} proceed diagrammatically: the four blocks $\Sigma^{ij}$ can be expanded
to higher order in $\hat V_R$ with generalized (cumulant-containing) diagrams,
and \MRGW{} is a selected resummation of the $\Sigma^{11}$ ring series.  The
\DHN{} organize everything in an exact Dyall--Hedin hierarchy with a screening
channel, a physical self-energy channel, and mixed/auxiliary channels, and
state which channel is truncated where; the next step is a two-frequency
kernel $W_D(\nu)\Gamma_D(\omega,\nu)$ and residual vertex corrections
$\delta W_D/\delta G_D$.  In the superoperator language of the \Pre{}, the
\MRGW{} bath of the \DHN{} can be read as an enlarged manifold of product
operators (an inactive or active charged operator times an MR-RPA excitation
operator) whose secondary block is taken diagonal, which is the usual
ADC-type reading of \GW{}~\cite{MarieTolleLoos2026}.

\paragraph{D8. Numerical evidence.}
\WFL{} provide the only numerical input: the Be satellite, the stretched-H$_2$
satellites, the O$_3$ state ordering, the orbital dependence (Table~S2), and,
most relevant here, the demonstration on H$_4$ that inserting bare
$\Sigma_1^{12/21}$ into Hall's equation breaks positivity (Fig.~6) while
changing peak positions only slightly.  Neither the \Pre{} nor the \DHN{}
contain numbers.  The H$_4$ observation is the empirical motivation for the
Hermitian embedding of S5, and the small shifts reported by \WFL{} suggest that
the mixed sector will matter mainly for satellite weights and for sum rules
(the first-moment error, ``First-moment error as a practical diagnostic'' in
the \DHN{}) rather than for principal ionization energies.

\paragraph{D9. Notational and definitional loose ends in the \Pre{}.}
(i) The subscripts $\mathrm{i,a}$ denote inactive/active in
$\bH_{\mathrm{i,a}}$, $\bSigma_{\mathrm{i,a}}$, and
$\ket{\Phi_{\mathrm i}}\otimes\ket{\Psi_{\mathrm a}}$, but core/virtual in the
manifold $\{a_i,a_a\}$; both \WFL{} and the \DHN{} use $P,Q$ for inactive and
$i,j$/$a,b$ for core/virtual.  (ii) $\tilde\bSigma_{\mathrm{i,i}}$ carries a
$\Sigma$ but is a difference of Green's functions.  (iii)
$\bSigma^{\text{corr}}$ is left undefined; the natural reading suggested by
the comparison is: its static and mixed first-order content is already in
$\bH_{\mathrm{i,a}}$, and what remains is the dynamic \MRGW{} bath (plus
screening of $\bH_{\mathrm{i,a}}$).  (iv) The killer condition is stated for
$\hat R^{N\pm1}_\mu$ but the metric for odd operators should be the
anticommutator; with $\hat R^{N+1}_\mu=\ket{\Psi^{N+1}_\mu}\bra{\Psi_0}$ both
are satisfied trivially.

\section{How the \Pre{} fits into the rest of the notes}
\label{sec:synthesis}

Three modifications turn the \Pre{} construction into the Objective~II Green's
function, and each one corresponds to a specific subsection of the \DHN{}.
\begin{enumerate}[leftmargin=2em]
\item \textbf{Retarded bookkeeping.}  Use signed pole energies
  $\kappa_{\mu\pm}=\pm(E^{N\pm1}_\mu-E^N_0)$ in $\bE$, the Dyall Fock energies
  in the inactive block, and write the coupling as a Hermitian matrix
  $[K_B]_{P\alpha}$, $[K_B]_{\alpha P}=[K_B]_{P\alpha}^*$ over \emph{all}
  inactive orbitals and \emph{all} active poles, including cross-sector pairs
  (core with $N+1$, virtual with $N-1$).  This is the content of
  ``Retarded spectral representation of the Dyall Green's function'' and
  ``Cross-sector pole pairs''.
\item \textbf{Screen the coupling.}  Replace $\bar v_{Pz,yw}$ in the
  three-operator part of $\bH_{\mathrm{i,a}}$, Eq.~\eqref{eq:H-ia-split}, by
  the parent-consistent static kernel
  $[\mathcal I_D^{\rm stat}]_{Pz,yw}=v^R_{Pz,yw}-[W_D(0)]_{Pw,yz}$
  (``A parent-consistent static kernel for the mixed Hall sector'').  The bare
  \Pre{} coupling is the $\mathcal I_D^{(0)}=\bar v_R$ member of the kernel
  hierarchy and is retained there as a control.
\item \textbf{Append the \MRGW{} bath.}  Border the matrix with the pole
  amplitudes $A$ and energies $E_{GW}$ of the spectral \MRGW{} self-energy,
  \begin{equation}
    \bbH
    =
    \begin{pmatrix}
      \bF_{\mathrm{i,i}}+[\Sigma_{\rm stat}]_{\mathrm{i,i}} & K_B^{\mathrm{i,a}} & A_{\mathrm i}\\
      K_B^{\mathrm{a,i}} & \bE & \bM^\dagger A_{\mathrm a}\\
      A_{\mathrm i}^\dagger & A_{\mathrm a}^\dagger\bM & E_{GW}
    \end{pmatrix},
    \qquad
    G(z)=\begin{pmatrix}\mathbf 1&\bM&0\end{pmatrix}(z-\bbH)^{-1}\begin{pmatrix}\mathbf 1&\bM&0\end{pmatrix}^{\dagger},
  \end{equation}
  which is $\bbH_{\text{mix-MR-}GW}$ of ``Common Hermitian matrix''.  Setting
  $K_B=0$ gives back \WFL{}'s \MRGW{} exactly (``Recovery of the existing
  \MRGW{} approximation''); setting $A=0$ and $W_D(0)\to v_R$ gives back the
  \Pre{}.
\end{enumerate}
A pilot implementation of this matrix, with a Davidson solver that follows
selected roots, consistency checks against full CI, and the H$_4$ example of
\WFL{}'s Fig.~6, is provided in the directory \texttt{numerics/} of the
repository (see its \texttt{README.md}).  With $K_B=0$ it reproduces the
\MRGW{} ionization energies of ozone reported in \WFL{}'s Table~S5 for
CAS(6,4) and CAS(8,5) to within $0.005$~eV, which confirms the identification
of their Dyson equation with the Hermitian bath representation.  With the
mixed couplings switched on, the state ordering $^2A_1<{}^2B_2<{}^2A_2$ is
preserved; the bare kernel changes the three ionization energies by at most
$0.2$~eV, the statically screened kernel shifts the $^2B_2$ and $^2A_2$ peaks
upward by $0.3$--$0.4$~eV.

Read in the other direction, the \Pre{} supplies two things the \DHN{} do not
have.  First, a derivation of the same-sector part of $K_B$ that is Hermitian
from the outset and needs no conversion from time-ordered Hall blocks, so that
the sign and prefactor questions of D1 and D6 can be settled by evaluating one
double commutator, Eq.~\eqref{eq:H-ia-explicit}; the cross-sector part has to
be taken from the first-order theory (D1).  Second, a transparent
physical reading of the higher-order terms of the $K_B$ resummation
(inactive--active configuration mixing between charged Dyall states, and the
induced inactive self-energy $\bSigma_{\mathrm{i,i}}$), which the \DHN{}
describe only as a ``chosen Hermitian resummation''.  The two bullets that
close the \Pre{} map onto the \WFL{} screening channels as follows: ``effects
within the inactive space itself'' is the core$\to$virtual ring channel, and
``couplings beyond the operator manifold'' are the three channels involving
active transition densities together with the dynamic
$(N\pm1)\otimes$(MR-RPA) bath.

\section{Points in the \Pre{} suggested for revision}
\label{sec:revisions}

\begin{enumerate}[leftmargin=2em]
\item Replace $\bF^{\text{HF}}_{\mathrm{i,i}}$ by the Dyall Fock operator with
  the active one-body density (D5), and state whether the core--virtual block
  $[\Sigma_1^{11}]_{aj}$ is kept (it vanishes only for CASSCF-optimized orbitals).
\item State the sign convention of the removal energies in $\bE^{N\pm1}$ and
  the conjugation pattern of the couplings (D6); correct $\bM\to\bM^\dagger$ in
  $\bSigma_{\mathrm{i,a}}$.
\item Remove the identification of $\tilde\bSigma_{\mathrm{i,i}}$ with
  $\Sigma^{22}$; describe $\bSigma_{\mathrm{i,i}}$ as the second-order
  inactive self-energy generated by the resummed inactive--active coupling (D3).
\item Sharpen the statement ``same structure as the initial Dyson equation in
  the \MRGW{} paper'' to: same block pattern as \WFL{}'s Eq.~(S20), but with the
  dressed $\bG_{\mathrm i}$ that removes the double poles of MR-GF1 (D4), i.e.\
  the \Pre{} construction is a Hermitian completion of the kind requested in
  \WFL{}'s Appendix~C.
\item Fix the $\mathrm{i,a}$ versus $i,a$ notation clash and rename
  $\tilde\bSigma_{\mathrm{i,i}}$ (D9).
\item Add the explicit form of $\bH_{\mathrm{i,a}}$, Eq.~\eqref{eq:H-ia-explicit},
  and note that it equals the same-sector part of
  $[T_D^\dagger\Sigma_{\rm stat}T_D+K_B^{(1)}]_{P\alpha}$ of Objective~II,
  while the cross-sector part is missing with the zeroth-order reference (D1);
  this closes the loop between the two parts of \texttt{mr-hedin.tex} and
  resolves the red-marked sign warning.
\item Define $\bSigma^{\text{corr}}$ operationally as the dynamic \MRGW{} bath
  plus the static screening of $\bH_{\mathrm{i,a}}$ (Sec.~\ref{sec:synthesis}).
\end{enumerate}

\begin{thebibliography}{99}

\bibitem{WFL2026}
Y.~Wang, W.-H.~Fang, and Z.~Li,
``A Diagrammatic Route to Multi-Reference $GW$ for Strongly Correlated Molecules'',
arXiv:2604.16013v2 [physics.chem-ph] (2026).

\bibitem{WFL2025}
Y.~Wang, W.-H.~Fang, and Z.~Li,
J.~Phys.~Chem.~Lett.\ \textbf{16}, 3047 (2025) [multireference RPA].

\bibitem{Hall1975}
A.~G.~Hall, J.~Phys.~A: Math.~Gen.\ \textbf{8}, 214 (1975).

\bibitem{Brouder2009}
C.~Brouder, in \emph{Quantum Field Theory: Competitive Models},
edited by B.~Fauser, J.~Tolksdorf, and E.~Zeidler (Birkh\"auser, Basel, 2009),
pp.~163--175; arXiv:0710.5652.

\bibitem{Dyall1995}
K.~G.~Dyall, J.~Chem.~Phys.\ \textbf{102}, 4909 (1995).

\bibitem{Hedin1965}
L.~Hedin, Phys.~Rev.\ \textbf{139}, A796 (1965).

\bibitem{MarieTolleLoos2026}
A.~Marie, J.~T\"olle, and P.-F.~Loos,
J.~Phys.~Chem.~Lett.\ \textbf{17}, 8541 (2026).

\bibitem{GoscinskiLukman1970}
O.~Goscinski and B.~Lukman, Chem.~Phys.~Lett.\ \textbf{7}, 573 (1970).

\bibitem{Ortiz2013}
J.~V.~Ortiz, WIREs Comput.~Mol.~Sci.\ \textbf{3}, 123 (2013).

\bibitem{Schirmer1983}
J.~Schirmer, L.~S.~Cederbaum, and O.~Walter, Phys.~Rev.~A \textbf{28}, 1237 (1983).

\bibitem{Sokolov2018}
A.~Y.~Sokolov, J.~Chem.~Phys.\ \textbf{149}, 204113 (2018).

\bibitem{ChatterjeeSokolov2019}
K.~Chatterjee and A.~Y.~Sokolov, J.~Chem.~Theory Comput.\ \textbf{15}, 5908 (2019).

\end{thebibliography}

\end{document}
